% !TeX program = pdflatex
\documentclass[11pt]{article}
\usepackage[a4paper,margin=1in]{geometry}
\usepackage[T1]{fontenc}
\usepackage{newtxtext}
\usepackage{amsmath,amssymb,amsthm}
\usepackage{newtxmath}
\usepackage{microtype}
\usepackage{mathtools}
\usepackage{booktabs}
\usepackage{multirow}
\usepackage{graphicx}
\usepackage{caption}
\usepackage{subcaption}
\usepackage{enumitem}
\usepackage{xcolor}
\usepackage{float}
\usepackage[section]{placeins}
\usepackage{bm}
\usepackage{hyperref}
\usepackage{tikz}
\usepackage{tikz-3dplot}
\usetikzlibrary{shapes,arrows,calc,decorations.pathreplacing,patterns}

% ---------- Hyperref setup ----------
\hypersetup{colorlinks=true,linkcolor=blue,citecolor=blue,urlcolor=blue}

% ---------- Cleveref ----------
\usepackage[capitalize,noabbrev]{cleveref}

% ---------- Theorem environments ----------
\renewcommand{\qedsymbol}{}
\newtheorem{theorem}{Theorem}[section]
\newtheorem{proposition}[theorem]{Proposition}
\newtheorem{lemma}[theorem]{Lemma}
\newtheorem{corollary}[theorem]{Corollary}
\theoremstyle{definition}
\newtheorem{definition}[theorem]{Definition}
\theoremstyle{remark}
\newtheorem{remark}[theorem]{Remark}

% ---------- Math helpers ----------
\newcommand{\E}{\mathbb{E}}
\newcommand{\R}{\mathbb{R}}
\newcommand{\KL}{\mathrm{KL}}
\newcommand{\softmax}{\mathrm{softmax}}
\DeclareMathOperator*{\argmax}{arg\,max}
\DeclareMathOperator*{\argmin}{arg\,min}

% Notations
\newcommand{\tauanc}{\tau_{\mathrm{anc}}}
\newcommand{\loss}{\mathcal{L}}
\newcommand{\target}{u_{\text{target}}}
\newcommand{\piref}{\pi_{\mathrm{ref}}}
\newcommand{\pik}{\pi_k}

\title{\textbf{W-GOPO: Weighted Group Orthogonalized Policy Optimization}\\
\large Escort-Weighted Advantages and Exact Conservation\\in the Hilbert Projection Framework}

\author{Wang Zixian\\
\small China Mobile Communications Group Shandong Co., Ltd. Tai'an Branch\\
\small \texttt{wangzixian@sd.chinamobile.com}}
\date{}

\begin{document}
\maketitle

\begin{abstract}
Orthogonalized Policy Optimization (OPO) and its group-level specialization GOPO reformulate LLM alignment as orthogonal projection in the Hilbert space $L^2(\pik)$, replacing the exponential KL geometry with constant-curvature $\chi^2$ optimization. However, these two algorithms occupy opposite ends of a fundamental trade-off. OPO operates in the full infinite-dimensional $L^2(\pik)$ with an $\alpha$-Escort sampling field that provides \emph{global peak-seeking}---updates driven by absolute reward quality---but introduces importance-sampling variance and requires solving for the chemical potential $\lambda^*$. GOPO achieves tractability by collapsing the projection onto finite group subspaces and leveraging the vanishing chemical potential under group-relative centering ($A_i = r_i - \bar{r}$), but this local centering \emph{amplifies credit misattribution} for uniformly poor batches and \emph{breaks probability conservation} when the Bounded Hilbert Projection truncation is active.

We present \textbf{Weighted Group Orthogonalized Policy Optimization (W-GOPO)}, which resolves this tension through three mechanisms: (i)~$\alpha$-Escort weighted advantages that restore absolute quality awareness within the group framework; (ii)~a GPU-parallel bisection solver for the exact chemical potential, guaranteeing strict probability conservation even under BHP truncation; and (iii)~a Pre-Squashing and Post-Projection pipeline that compresses raw rewards through a bounded $\tanh$ surrogate before zero-mean projection, preventing reward outliers from destabilizing the Hilbert geometry while preserving exact sparsity. W-GOPO is a strict generalization of GOPO---setting $\alpha = 0$ and $c \to \infty$ recovers the original algorithm---while inheriting OPO's global sensitivity to reward quality. Experiments on mathematical reasoning benchmarks demonstrate that W-GOPO achieves the highest out-of-distribution generalization among all tested methods, with monotonically improving validation accuracy and the healthiest gradient dynamics.
\end{abstract}

%====================================================================%
\section{Introduction}
\label{sec:intro}

Aligning Large Language Models (LLMs) with human preferences via Reinforcement Learning from Human Feedback (RLHF)~\cite{Christiano2017RLHF,Ouyang2022InstructGPT} has become a standard practice. The dominant paradigm employs algorithms such as PPO~\cite{Schulman2017PPO}, DPO~\cite{Rafailov2023DPO}, and GRPO~\cite{Shao2024GRPO}, all of which share a common reliance on Kullback--Leibler (KL) divergence for trust-region enforcement. As demonstrated in prior work~\cite{Wang2025OPO,Wang2025GOPO}, the exponential geometry induced by KL divergence leads to gradient saturation in high-confidence regimes, motivating a shift to the Hilbert function space $L^2(\pik)$.

The Hilbert space reformulation, first proposed in OPO~\cite{Wang2025OPO} and specialized to group sampling in GOPO~\cite{Wang2025GOPO}, transforms the nonlinear probability simplex constraint into a linear orthogonality condition $\langle v, \mathbf{1}\rangle_{\pik} = 0$, where $v(y) = \pi(y)/\pik(y) - 1$ is the density fluctuation field. The optimal policy update is given by the Hilbert Projection Theorem, yielding constant Hessian curvature, non-saturating gradients, and exact sparsity via the Bounded Hilbert Projection (BHP). These properties have been shown to sustain learning in regimes where clipping-based methods plateau~\cite{Wang2025GOPO}.

\paragraph{The OPO--GOPO Trade-off.}
Despite sharing the same Hilbert-space foundation, OPO and GOPO make qualitatively different engineering trade-offs that lead to distinct failure modes:

\begin{itemize}[leftmargin=2em]
    \item \textbf{OPO}~\cite{Wang2025OPO} operates in the full $L^2(\pik)$ with an $\alpha$-Escort sampling field (unnormalized weight $\tilde{\omega}_\alpha$, subsequently centered) that provides \emph{global peak-seeking}: the update magnitude is controlled by the \emph{absolute} quality of generated responses. When a batch consists entirely of poor outputs, all $|A_i|$ are small in absolute terms, and the escort weight naturally suppresses the entire update. However, OPO requires solving for the chemical potential $\lambda^*$ in the infinite-dimensional space and introduces importance-sampling variance.
    
    \item \textbf{GOPO}~\cite{Wang2025GOPO} collapses the projection onto finite group subspaces. By using group-relative centering ($A_i = r_i - \bar{r}$, guaranteeing $\sum A_i = 0$), GOPO eliminates the chemical potential entirely ($\lambda^* = 0$ in the BHP-inactive regime), yielding an unconstrained quadratic loss of remarkable simplicity. However, this local centering creates two structural pathologies that we identify in this work.
\end{itemize}

\paragraph{Problem 1: Reward Amplification in Degenerate Batches.}
Consider a batch where all $G$ sampled responses are of uniformly poor quality. Under GOPO's group-relative centering ($A_i = r_i - \bar{r}$), the \emph{least-bad} response necessarily receives a positive advantage whose magnitude depends only on the relative spread within the group---not on the absolute reward level. Because the sentence-level policy gradient broadcasts this advantage to every token in the response, all tokens---including low-quality ones---receive positive gradient signals. The model is forced to reinforce mediocre behavior. In OPO's framework, the escort weight $\tilde{\omega}_\alpha$ would suppress this entire batch (all absolute advantages are small); GOPO's centering cannot distinguish ``uniformly poor with small spread'' from ``uniformly excellent with similar spread,'' treating both identically.

\paragraph{Problem 2: Conservation Breakdown under BHP Truncation.}
GOPO's central structural simplification---the vanishing chemical potential $\lambda^* = 0$---relies on the linearity of the zero-mean condition $\sum A_i = 0$. When the BHP constraint $v(y) \geq -1$ is active, the $\max(-1, \cdot)$ truncation operator breaks this linearity. Specifically, if some response has $A_k/\mu < -1$, the truncation clamps $v_k^*$ to $-1$ instead of the more negative value $A_k/\mu$, creating a positive residual:
\begin{equation}
\label{eq:conservation_break_intro}
    \sum_{i=1}^G v_i^* = \sum_{i=1}^G \max\!\left(-1,\; \frac{A_i}{\mu}\right) > 0
\end{equation}
This ``phantom probability''---probability mass created from nowhere---violates the geometric foundation of the entire Hilbert projection framework.

\paragraph{This Work.}
We present \textbf{Weighted Group Orthogonalized Policy Optimization (W-GOPO)}, which resolves both pathologies while preserving GOPO's computational simplicity. Our approach combines three complementary mechanisms:

\begin{enumerate}[leftmargin=2em]
    \item \textbf{$\alpha$-Escort weighted advantages} (\cref{subsec:escort}): We integrate OPO's escort sampling field into the group framework, providing absolute quality awareness that superlinearly suppresses updates for degenerate batches while focusing on truly excellent responses.
    \item \textbf{Exact chemical potential via GPU-parallel bisection} (\cref{subsec:bisection}): We replace the $\lambda^* = 0$ assumption with an exact solver for the nonlinear conservation equation, restoring strict probability conservation under BHP truncation at negligible computational cost.
    \item \textbf{Pre-Squashing and Post-Projection} (\cref{subsec:squash}): We introduce a bounded $\tanh$ surrogate reward that compresses raw rewards before zero-mean projection, preventing reward outliers from destabilizing the Hilbert geometry while preserving BHP's exact sparsity guarantee.
\end{enumerate}

W-GOPO is a strict generalization of GOPO: setting $\alpha = 0$ (no escort weighting) and $c \to \infty$ (no squashing) recovers the original algorithm exactly. The resulting method inherits GOPO's computational simplicity---no critic network, no importance-sampling variance---while recovering OPO's global sensitivity to absolute reward quality.

\paragraph{Contributions.}
\begin{itemize}[leftmargin=2em]
    \item We rigorously characterize two structural pathologies of GOPO: reward amplification in degenerate batches (\cref{subsec:amplification}) and probability conservation breakdown under BHP truncation (\cref{subsec:bhp_break}).
    \item We derive W-GOPO as the minimal extension of GOPO that addresses both pathologies, proving that it subsumes GOPO as a special case (\cref{thm:reduction}).
    \item We provide a complete GPU-parallel implementation of the exact chemical potential solver with convergence guarantees (\cref{lem:monotonicity}, \cref{alg:bisection}).
    \item Experiments on mathematical reasoning benchmarks demonstrate that W-GOPO achieves the highest out-of-distribution generalization (49\% on MATH Level~4), surpassing OPO (48\%), GOPO (47\%), and clipping-based methods (44\%), with the healthiest gradient dynamics and entropy preservation.
\end{itemize}

%====================================================================%
\section{Preliminaries: The Hilbert Projection Framework}
\label{sec:prelim}

We review the shared mathematical foundation of OPO~\cite{Wang2025OPO} and GOPO~\cite{Wang2025GOPO}. For a given prompt $x$, let $\pik(\cdot|x)$ denote the current reference policy. The Hilbert space $\mathcal{H} = L^2(\pik)$ is equipped with inner product $\langle f, g \rangle_{\pik} = \E_{\pik}[f(y)g(y)]$.

\subsection{Density Fluctuation and Probability Conservation}

The density fluctuation field $v(y) = \pi(y)/\pik(y) - 1$ transforms the nonlinear probability simplex constraint into a linear orthogonality condition:
\begin{equation}
\label{eq:conservation}
    \sum_y \pi(y) = 1 \quad \Longleftrightarrow \quad \langle v, \mathbf{1}\rangle_{\pik} = \E_{\pik}[v] = 0
\end{equation}
This defines the closed subspace $\mathcal{H}_0 = \{f \in \mathcal{H} : \langle f, \mathbf{1}\rangle_{\pik} = 0\}$ of codimension one.

\subsection{The Work-Dissipation Functional}

Given a driving field $g(y)$ (derived from advantages), the alignment objective maximizes external work against quadratic dissipation:
\begin{equation}
\label{eq:work_diss}
    \max_{v \in \mathcal{H}_0}\; \mathcal{J}(v) = \langle g, v \rangle_{\pik} - \frac{\mu}{2}\|v\|^2_{\pik}
\end{equation}
As shown in~\cite{Wang2025GOPO}, this functional is algebraically equivalent to minimizing the squared $L^2$ distance $\frac{\mu}{2}\|v - g/\mu\|^2_{\pik}$ between the valid fluctuation $v \in \mathcal{H}_0$ and the unconstrained target $u^* = g/\mu$.

\subsection{The Orthogonal Projection Solution}

By the Hilbert Projection Theorem~\cite{Wang2025OPO}, the unique maximizer is:
\begin{equation}
\label{eq:projection}
    v^* = P_{\mathcal{H}_0}(u^*) = \frac{1}{\mu}\big(g - \underbrace{\E_{\pik}[g]}_{\lambda^*} \cdot \mathbf{1}\big)
\end{equation}
The subtracted term $\lambda^* = \E_{\pik}[g]$ is the \emph{chemical potential}: the Lagrange multiplier enforcing probability conservation, so that $v^* = (g - \lambda^*)/\mu$. Geometrically, $\lambda^*$ is the component of the driving field $g$ along the illegal direction $\mathbf{1}$ that is sliced off by the projection.

\subsection{Bounded Hilbert Projection (BHP)}

Enforcing non-negativity $\pi(y) \geq 0$, equivalently $v(y) \geq -1$, upgrades the projection to the \emph{Bounded Hilbert Projection}~\cite{Wang2025GOPO}:
\begin{equation}
\label{eq:bhp}
    v^*(y) = \max\!\left(-1,\; \frac{g(y) - \lambda^*}{\mu}\right)
\end{equation}
where $\lambda^*$ now satisfies the nonlinear equation $\E_{\pik}[v^*] = 0$. Actions with $g(y) < \lambda^* - \mu$ receive $v^*(y) = -1$, i.e., target probability \emph{exactly zero}---the \emph{exact sparsity} property.

\subsection{GOPO: The Group-Level Specialization}

GOPO transitions the projection from $L^2(\pik)$ to a finite-dimensional empirical subspace induced by group sampling. For a group $\mathcal{G} = \{y_1, \ldots, y_G\}$ with rewards $r_i$, group-relative advantages $A_i = r_i - \bar{r}$ satisfy $\sum A_i = 0$ by construction. Substituting $g(y_i) = A_i$ into \cref{eq:projection}:
\begin{equation}
    \lambda^* = \frac{1}{G}\sum_{i=1}^G A_i = 0
\end{equation}
The chemical potential \emph{vanishes}, collapsing the constrained projection into the unconstrained GOPO loss:
\begin{equation}
\label{eq:gopo_loss}
    \loss_{\text{GOPO}}(\theta) = -\frac{1}{G}\sum_{i=1}^G \left[ A_i \rho_\theta(y_i) - \frac{\mu}{2}\big(\rho_\theta(y_i) - 1\big)^2 \right]
\end{equation}
where $\rho_\theta(y_i) = \pi_\theta(y_i|x)/\pik(y_i|x)$ is the density ratio.

%====================================================================%
\section{Structural Limitations of GOPO}
\label{sec:problems}

We now precisely characterize the two pathologies that arise from GOPO's structural simplifications.

\subsection{Pathology 1: Reward Amplification in Degenerate Batches}
\label{subsec:amplification}

The group-relative centering $A_i = r_i - \bar{r}$ maps raw rewards to a zero-mean representation with scale determined by the group's internal variance, irrespective of absolute reward quality.

\begin{proposition}[Amplification Bound]
\label{prop:amplification}
Let $r_{\max} = \max_i r_i$ and $\bar{r} = \frac{1}{G}\sum_i r_i$. Under group-relative centering, the maximum positive advantage is:
\begin{equation}
    A_{\max} = r_{\max} - \bar{r}
\end{equation}
This quantity depends only on the \emph{relative spread} of rewards within the group, not on their absolute magnitude. For example, a uniformly poor group with rewards $\{0.02, 0.05, 0.08, 0.03, 0.04, 0.06\}$ (mean $0.047$, spread $0.06$) produces advantages of \emph{identical} scale to a uniformly excellent group $\{0.62, 0.65, 0.68, 0.63, 0.64, 0.66\}$ (mean $0.647$, spread $0.06$). Group-relative centering cannot distinguish between these two qualitatively different situations.
\end{proposition}

The severity of this pathology is compounded by the sentence-level credit assignment structure of autoregressive language models. When the policy gradient $\nabla_\theta \log\pi_\theta(y)$ distributes the advantage signal uniformly across all tokens in $y$, inflated advantages for mediocre responses force the model to reinforce every token---including low-quality ones---with positive gradient signals.

\begin{remark}[Contrast with OPO's Escort Field]
\label{rem:escort_contrast}
In OPO~\cite{Wang2025OPO}, the unnormalized escort weight $\tilde{\omega}_\alpha$ is centered and then used as the driving field. For a degenerate batch with mean absolute advantage $\bar{A} \to 0$, the exponential term $A_i \exp(A_i) \approx A_i + A_i^2$ introduces a superlinear $O(\bar{A}^2)$ component that suppresses the total update magnitude. GOPO's group-relative centering, by contrast, produces advantages whose scale depends only on the intra-group spread, providing updates of the same magnitude for uniformly poor batches as for high-quality ones with similar spread.
\end{remark}

\subsection{Pathology 2: Probability Conservation Breakdown under BHP}
\label{subsec:bhp_break}

The vanishing chemical potential $\lambda^* = 0$ in GOPO is a consequence of the linear identity $\sum A_i = 0$. When the BHP truncation $v(y) \geq -1$ is active, the $\max(-1, \cdot)$ operator introduces a nonlinearity that invalidates this identity.

\begin{proposition}[Conservation Breakdown]
\label{prop:conservation_break}
Let $S = \{i : A_i/\mu < -1\}$ denote the set of indices where BHP truncation is active, with $|S| = k > 0$. Under the GOPO assumption $\lambda^* = 0$, the empirical conservation residual is:
\begin{equation}
    \Delta \equiv \sum_{i=1}^G v_i^* = \sum_{i=1}^G \max\!\left(-1,\; \frac{A_i}{\mu}\right) = -\sum_{i \in S}\left(\frac{A_i}{\mu} + 1\right) > 0
\end{equation}
The excess $\Delta > 0$ represents ``phantom probability'' injected by the truncation. Each clamped index $i \in S$ contributes $|A_i/\mu + 1|$ units of excess mass.
\end{proposition}

\begin{proof}
Decomposing the sum over truncated and non-truncated indices:
\begin{align}
    \sum_{i=1}^G v_i^* &= \sum_{i \notin S} \frac{A_i}{\mu} + \sum_{i \in S} (-1) 
    = \frac{1}{\mu}\sum_{i=1}^G A_i - \frac{1}{\mu}\sum_{i \in S} A_i - k
    = 0 - \sum_{i \in S}\!\left(\frac{A_i}{\mu} + 1\right)
\end{align}
Since $A_i/\mu < -1$ for $i \in S$, each summand $A_i/\mu + 1 < 0$, so the total is positive.
\end{proof}

\begin{remark}[Physical Consequence]
The excess $\Delta > 0$ means the target policy $\pi_{\text{target}}(y_i) = \pik(y_i)(1 + v_i^*)$ no longer sums to one over the group. The policy is being steered toward a distribution that allocates \emph{more total probability} than available---a physically impossible target. While the MSE regression in GOPO may partially absorb this discrepancy in practice, the geometric foundation of the Hilbert projection is formally violated.
\end{remark}

%====================================================================%
\section{W-GOPO: Resolving the OPO--GOPO Tension}
\label{sec:wgopo}

We now present W-GOPO, designed as the \emph{minimal extension} of GOPO that addresses both pathologies while preserving its core structural properties: constant curvature, non-saturating gradients, and no critic network.

\subsection{Mechanism 1: $\alpha$-Escort Weighted Advantages}
\label{subsec:escort}

W-GOPO replaces GOPO's raw advantages with an escort-weighted driving signal. Following OPO's principle of using unnormalized weights that are subsequently centered~\cite{Wang2025OPO}, we define the \emph{interpolated escort weight} for each response $y_i$ with advantage $A_i$:
\begin{equation}
\label{eq:escort}
    \tilde{\omega}_\alpha(A_i) = (1 - \alpha)\, A_i \;+\; \alpha\, A_i \cdot \exp(A_i), \quad \alpha \in [0, 1]
\end{equation}
This definition is designed so that $\alpha$ smoothly interpolates between GOPO and OPO-style peak-seeking:
\begin{itemize}[leftmargin=2em]
    \item $\alpha = 0$: $\tilde{\omega}_0(A_i) = A_i$, recovering standard GOPO exactly.
    \item $\alpha > 0$: The $A_i \exp(A_i)$ term exponentially amplifies high-advantage responses, providing global peak-seeking.
    \item $\alpha = 1$: $\tilde{\omega}_1(A_i) = A_i \exp(A_i)$, the strongest peak-seeking in the family.
\end{itemize}
When Pre-Squashing ensures $|A_i| \leq 2c$ (\cref{subsec:squash}), the weight is naturally bounded:
\begin{equation}
    |\tilde{\omega}_\alpha(A_i)| \leq (1 - \alpha)\,|A_i| + \alpha\,|A_i|\,e^{|A_i|} \leq 2c \cdot \big[(1-\alpha) + \alpha\,e^{2c}\big]
\end{equation}
eliminating the numerical instability (inf/underflow) associated with power-law weights $|A|^{p}$ when $p = 1/(1-\alpha) \to \infty$.

Because $\tilde{\omega}_\alpha$ is nonlinear for $\alpha > 0$, the zero-mean property is broken: $\sum \tilde{\omega}_\alpha(A_i) \neq 0$ in general. The driving field must therefore be explicitly centered to define the chemical potential:
\begin{equation}
\label{eq:escort_center}
    g_i = \tilde{\omega}_\alpha(A_i) - \bar{\omega}_\alpha, \qquad \bar{\omega}_\alpha = \frac{1}{G}\sum_{j=1}^G \tilde{\omega}_\alpha(A_j)
\end{equation}
This centering is the exact group-level analog of OPO's conservation projection $\mathcal{P}_{\mathcal{H}_0}$~\cite{Wang2025OPO}: the centered field $g_i$ satisfies $\sum g_i = 0$, and $\bar{\omega}_\alpha$ plays the role of the chemical potential in the BHP-inactive regime ($\lambda^* = 0$ when applied to $g$).

\begin{theorem}[Gradient Suppression for Degenerate Batches]
\label{thm:suppression}
Let $\bar{A} = \frac{1}{G}\sum_i |A_i|$ denote the mean absolute advantage. For $\alpha > 0$, the total escort-weighted update magnitude satisfies:
\begin{equation}
    \sum_{i=1}^G |\tilde{\omega}_\alpha(A_i)| \leq G\bar{A}\big[(1-\alpha) + \alpha\,e^{\bar{A}}\big] \cdot \frac{\sum|A_i|e^{|A_i|}}{G\bar{A} e^{\bar{A}}}
\end{equation}
By Jensen's inequality (convexity of $x \mapsto x e^x$ on $[0,\infty)$), the ratio on the right is $\geq 1$, and the dominant behavior for small $\bar{A}$ is $O(\bar{A})$ at $\alpha = 0$ versus $O(\bar{A}^2)$ at $\alpha > 0$ (since $e^{\bar{A}} \approx 1 + \bar{A}$ for small $\bar{A}$, making $\bar{A} e^{\bar{A}} \approx \bar{A} + \bar{A}^2$). This provides superlinear suppression for degenerate batches ($\bar{A} \to 0$) while amplifying genuinely excellent responses.
\end{theorem}

When BHP truncation is active, the centered field $g_i$ is used to compute the BHP target, and $\lambda^*$ is solved via bisection to enforce exact conservation.

\subsection{Mechanism 2: Exact Chemical Potential via Bisection}
\label{subsec:bisection}

With the centered driving field $g_i$ from~\eqref{eq:escort_center}, the BHP solution becomes:
\begin{equation}
\label{eq:bhp_wgopo}
    v_i^*(\lambda) = \max\!\left(-1,\; \frac{g_i - \lambda}{\mu}\right)
\end{equation}
and $\lambda^*$ must satisfy the conservation equation:
\begin{equation}
\label{eq:conservation_eq}
    f(\lambda) \equiv \sum_{i=1}^G \max\!\left(-1,\; \frac{g_i - \lambda}{\mu}\right) = 0
\end{equation}

\begin{lemma}[Monotonicity and Unique Root]
\label{lem:monotonicity}
The function $f(\lambda)$ is continuous and monotonically non-increasing on $\R$, with $\lim_{\lambda \to -\infty} f(\lambda) = +\infty$ and $\lim_{\lambda \to +\infty} f(\lambda) = -G$. Furthermore, $f$ is \emph{strictly} decreasing on the interval $\{\lambda : f(\lambda) > -G\}$ (i.e., wherever at least one summand is unclamped). Since $f(-\infty) = +\infty > 0 > -G = f(+\infty)$, the unique root $\lambda^*$ satisfying $f(\lambda^*) = 0$ exists and lies in the strictly decreasing region.
\end{lemma}

\begin{proof}
Each summand $\max(-1, (g_i - \lambda)/\mu)$ is non-increasing in $\lambda$: it decreases linearly with slope $-1/\mu$ when $(g_i - \lambda)/\mu > -1$ (unclamped), and is constant at $-1$ when clamped. Their sum $f$ is therefore non-increasing. Strict decrease holds on $\{\lambda : f(\lambda) > -G\}$, because $f > -G$ implies at least one summand is unclamped and contributes slope $-1/\mu < 0$. The limit properties follow directly: as $\lambda \to -\infty$, all terms are unclamped and $(g_i - \lambda)/\mu \to +\infty$; as $\lambda \to +\infty$, all terms clamp to $-1$, giving $f = -G$. Since $0 \in (-G, +\infty)$, the root $\lambda^*$ lies in the strictly decreasing region and is unique.
\end{proof}

The monotonicity of $f$ guarantees that bisection search converges to $\lambda^*$ with precision $\epsilon$ in $O(\log(1/\epsilon))$ iterations. We implement this as a fully batched GPU operation:

\begin{figure}[htbp]
\centering
\fbox{\parbox{0.92\textwidth}{
\textbf{Algorithm 1: GPU-Parallel Bisection for the Chemical Potential $\lambda^*$}
\vspace{0.3em}
\hrule
\vspace{0.3em}
\textbf{Input:} $\bm{g} \in \R^{B \times G}$ (centered escort field from~\eqref{eq:escort_center}), $\mu > 0$, tolerance $\epsilon$, max iterations $T$

\textbf{Initialize:}
$\bm{\lambda}_{\text{high}} \leftarrow \max_i g_i + \mu$ (per-batch), \quad
$\bm{\lambda}_{\text{low}} \leftarrow \min_i g_i - \mu G$ (per-batch)

\textbf{for} $t = 1, \ldots, T$ \textbf{do}
\begin{enumerate}[leftmargin=2em, itemsep=0pt]
    \item $\bm{\lambda}_{\text{mid}} \leftarrow (\bm{\lambda}_{\text{low}} + \bm{\lambda}_{\text{high}}) / 2$
    \item $\bm{v} \leftarrow \text{clamp}((\bm{g} - \bm{\lambda}_{\text{mid}}) / \mu, \; \min\!=\!-1)$ \hfill \textit{(BHP truncation)}
    \item $\bm{s} \leftarrow \sum_{i=1}^G v_i$ \hfill \textit{(conservation residual, per-batch)}
    \item $\bm{\lambda}_{\text{low}} \leftarrow \text{where}(\bm{s} > 0, \;\bm{\lambda}_{\text{mid}},\; \bm{\lambda}_{\text{low}})$
    \item $\bm{\lambda}_{\text{high}} \leftarrow \text{where}(\bm{s} \leq 0, \;\bm{\lambda}_{\text{mid}},\; \bm{\lambda}_{\text{high}})$
    \item \textbf{if} $\max(\bm{\lambda}_{\text{high}} - \bm{\lambda}_{\text{low}}) < \epsilon$: \textbf{break}
\end{enumerate}
\textbf{Output:} $\bm{\lambda}^* \leftarrow (\bm{\lambda}_{\text{low}} + \bm{\lambda}_{\text{high}}) / 2$, \;
$\bm{v}^* \leftarrow \text{clamp}((\bm{g} - \bm{\lambda}^*) / \mu, \;\min\!=\!-1)$
}}
\caption{Bisection search for the exact chemical potential. All operations are batched over $B$ prompts on GPU. The \texttt{clamp} + \texttt{sum} + \texttt{where} operations run on $(B, G)$ tensors with $T \leq 40$ and $G \leq 16$, adding $< 1$\,ms overhead per training step.}
\label{alg:bisection}
\end{figure}

\begin{proposition}[Convergence Rate]
\label{prop:convergence}
\cref{alg:bisection} reduces the bracket width $|\lambda_{\text{high}} - \lambda_{\text{low}}|$ below $\epsilon$ in at most $T = \lceil\log_2((g_{\max} - g_{\min} + \mu(G+1))/\epsilon)\rceil$ iterations. Here $\epsilon$ denotes the $\lambda$-precision (bracket width); the conservation residual $|f(\lambda)|$ at the midpoint satisfies $|f(\lambda)| \leq (G/\mu) \cdot \epsilon / 2$ upon termination, since $f$ has Lipschitz constant at most $G/\mu$ on the unclamped region.
\end{proposition}

\subsection{Mechanism 3: Pre-Squashing and Post-Projection}
\label{subsec:squash}

Even with escort weighting, the raw advantage can exhibit extreme outliers due to reward model miscalibration or inherent task difficulty variation. Any direct nonlinear transformation applied \emph{after} the zero-mean projection would destroy the $\sum A_i = 0$ property, causing uncontrolled resurrection of the chemical potential~\cite{Wang2025OPO}. W-GOPO resolves this via a two-stage pipeline that applies the nonlinearity \emph{before} projection:

\paragraph{Stage 1: Pre-Squashing.}
Raw rewards $r_i$ are centered and compressed through a bounded $\tanh$ surrogate:
\begin{equation}
\label{eq:squash}
    \tilde{u}_i = c \cdot \tanh\!\left(\frac{r_i - \bar{r}}{c}\right)
\end{equation}
where $c > 0$ controls the saturation width and $\bar{r} = \frac{1}{G}\sum_j r_j$.

\paragraph{Stage 2: Post-Projection.}
The squashed values are explicitly zero-mean projected:
\begin{equation}
\label{eq:post_proj}
    A_i = \tilde{u}_i - \frac{1}{G}\sum_{j=1}^G \tilde{u}_j
\end{equation}

\begin{proposition}[Properties of the Squashed Advantage]
\label{prop:squash}
The Pre-Squashing pipeline satisfies:
\begin{enumerate}[leftmargin=2em]
    \item \textbf{Boundedness:} $|A_i| \leq 2c$ for all $r_i$, providing a hard cap on the driving signal.
    \item \textbf{Local linearity:} $\tilde{u}_i \approx r_i - \bar{r}$ when $|r_i - \bar{r}| \ll c$, preserving local reward structure.
    \item \textbf{Conservation:} $\sum A_i = 0$ by construction. When no BHP truncation is active, the chemical potential vanishes exactly. When truncation is active, the bisection algorithm restores conservation.
    \item \textbf{1-Lipschitz:} $|\partial \tilde{u}_i / \partial r_i| \leq 1$, ensuring stable gradient flow.
\end{enumerate}
\end{proposition}

The scale parameter $c$ provides a principled defense against reward hacking (Goodhart's Law~\cite{Gao2023Scaling}): regardless of how extreme the raw reward becomes, the effective advantage is bounded by $2c$, and the escort-weighted signal by $2c[(1-\alpha) + \alpha e^{2c}]$---a finite, controllable quantity.

%====================================================================%
\section{The Complete W-GOPO Algorithm}
\label{sec:algorithm}

\subsection{Loss Function}

Combining all three mechanisms, the W-GOPO loss is:
\begin{equation}
\label{eq:wgopo_loss}
\boxed{
    \loss_{\text{W-GOPO}}(\theta) = \frac{1}{G}\sum_{i=1}^G \left( \rho_\theta(y_i) - (1 + v_i^*) \right)^2
}
\end{equation}
where $v_i^*$ is the BHP target computed from the escort-weighted, squashed advantages via \cref{alg:bisection}. The target $1 + v_i^*$ is treated as a \emph{detached pseudo-label}: it is computed from the current batch's rewards and advantages but does not propagate gradients.

\begin{remark}[MSE as Hilbert Distance]
The loss~\eqref{eq:wgopo_loss} is the parameterized version of the Hilbert Projection distance $\|v - v^*\|^2$: the policy ratio $\rho_\theta = 1 + v_\theta$ is steered toward the nearest feasible point $1 + v^*$ in $\mathcal{K} = \mathcal{H}_0 \cap \mathcal{C}_+$. The MSE form inherits constant Hessian $2I/G$ from the quadratic structure, ensuring non-saturating gradients.
\end{remark}

\subsection{Algorithm Summary}

\begin{figure}[htbp]
\centering
\fbox{\parbox{0.92\textwidth}{
\textbf{Algorithm 2: Weighted Group Orthogonalized Policy Optimization (W-GOPO)}
\vspace{0.3em}
\hrule
\vspace{0.3em}
\textbf{Input:} Policy $\pi_\theta$, stiffness $\mu$, escort exponent $\alpha$, squash scale $c$, group size $G$

\textbf{for} each iteration \textbf{do}
\begin{enumerate}[leftmargin=2em, itemsep=0pt]
    \item \textbf{Anchor:} $\pik \leftarrow \pi_\theta$ \hfill \textit{(on-policy anchoring)}
    \item \textbf{Sample:} For each prompt $x$, generate $\{y_1, \ldots, y_G\} \sim \pik(\cdot|x)$
    \item \textbf{Score:} Compute raw rewards $r_i = R(x, y_i)$
    \item \textbf{Pre-Squash:} $\tilde{u}_i = c\tanh((r_i - \bar{r})/c)$ \hfill \textit{(bounded surrogate)}
    \item \textbf{Post-Project:} $A_i = \tilde{u}_i - \frac{1}{G}\sum_j \tilde{u}_j$ \hfill \textit{($\sum A_i = 0$)}
    \item \textbf{Escort Weight:} $\tilde{\omega}_i = (1{-}\alpha)\,A_i + \alpha\,A_i\exp(A_i)$ \hfill \textit{(Eq.~\ref{eq:escort})}
    \item \textbf{Center:} $g_i = \tilde{\omega}_i - \frac{1}{G}\sum_j \tilde{\omega}_j$ \hfill \textit{($\sum g_i = 0$, conservation projection)}
    \item \textbf{Solve $\lambda^*$:} Bisection on $f(\lambda) = \sum \max(-1, (g_i - \lambda)/\mu) = 0$ \hfill \textit{(Alg.~1)}
    \item \textbf{BHP Target:} $v_i^* = \max(-1, (g_i - \lambda^*)/\mu)$ \hfill \textit{(exact sparsity)}
    \item \textbf{Compute Ratios:} $\rho_\theta(y_i) = \pi_\theta(y_i|x)/\pik(y_i|x)$
    \item \textbf{Loss:} $\loss = \frac{1}{G}\sum_i (\rho_\theta(y_i) - (1 + v_i^*))^2$ \hfill \textit{(Hilbert distance)}
    \item \textbf{Update:} $\theta \leftarrow \theta - \eta\nabla_\theta \loss$
\end{enumerate}
\textbf{end for}
}}
\caption{W-GOPO algorithm. Steps 4--5 implement Pre-Squashing and Post-Projection for reward robustness. Step 6 computes the interpolated escort weight and Step 7 centers it (conservation projection). Steps 8--9 solve for the exact BHP target with strict probability conservation. Steps 10--11 implement the Hilbert distance regression.}
\label{alg:wgopo}
\end{figure}

\subsection{Reduction to Prior Methods}

\begin{theorem}[GOPO as a Special Case]
\label{thm:reduction}
Setting $\alpha = 0$ and $c \to \infty$ in W-GOPO recovers the standard GOPO loss~\cite{Wang2025GOPO} in the BHP-inactive regime (i.e., when $A_i/\mu \geq -1$ for all $i$, so that no truncation occurs). When BHP truncation is active, W-GOPO with $\alpha = 0$, $c \to \infty$ differs from GOPO precisely by solving for the correct $\lambda^* \neq 0$---this is by design, as GOPO's $\lambda^* = 0$ assumption violates conservation in this regime (\cref{prop:conservation_break}).
\end{theorem}

\begin{proof}
When $\alpha = 0$: $\tilde{\omega}_0(A_i) = (1-0)A_i + 0 \cdot A_i\exp(A_i) = A_i$, so the centered field $g_i = A_i - \bar{A} = A_i$ (since $\sum A_i = 0$ after Post-Projection). When $c \to \infty$: $\tilde{u}_i = c\tanh((r_i - \bar{r})/c) \to r_i - \bar{r}$. In the BHP-inactive regime ($A_i/\mu \geq -1$ for all $i$), the bisection solver yields $\lambda^* = 0$ (since $\sum \max(-1, A_i/\mu) = \sum A_i/\mu = 0$), and the loss becomes $\frac{1}{G}\sum_i (\rho_i - 1 - A_i/\mu)^2$, equivalent to GOPO's quadratic form~\eqref{eq:gopo_loss}.
\end{proof}

%====================================================================%
\section{Theoretical Analysis}
\label{sec:analysis}

\subsection{Constant Curvature Preservation}

\begin{theorem}[Geometry Preservation]
\label{thm:geometry}
The W-GOPO loss inherits the following structural properties from the Hilbert projection framework:
\begin{enumerate}[leftmargin=2em]
    \item \textbf{Constant Hessian:} $\nabla^2_\rho \loss = 2I/G$, independent of advantages, rewards, or policy state.
    \item \textbf{Non-saturating gradients:} $\nabla_\rho \loss = 2(\rho - (1+v^*))/G$ is linear in the displacement from target, with magnitude proportional to distance from equilibrium.
    \item \textbf{Exact sparsity:} The BHP mechanism with bisection-solved $\lambda^*$ assigns $v_i^* = -1$ (target probability zero) for responses with $g_i < \lambda^* - \mu$.
\end{enumerate}
\end{theorem}

The escort weighting and Pre-Squashing modify \emph{which equilibrium} the policy converges to (by reshaping the target $v_i^*$), but not the \emph{convergence dynamics} (which are governed by the constant Hessian).

\subsection{Escort Weighting as Credit Misattribution Control}

\begin{corollary}[Contrast with GOPO]
\label{cor:contrast}
Under GOPO's group-relative centering ($\alpha = 0$), the advantage scale depends only on the intra-group reward spread, providing updates of the same magnitude for uniformly poor batches as for high-quality ones with similar spread. Under W-GOPO with $\alpha > 0$, the $A_i \exp(A_i)$ component introduces superlinear suppression: for small $\bar{A}$, $|\tilde{\omega}_\alpha(A)| \approx |A| + \alpha A^2$, so the total update acquires an $O(\bar{A}^2)$ damping term that vanishes for degenerate batches. Higher $\alpha$ amplifies this effect.
\end{corollary}

\subsection{Robustness Bound}

\begin{proposition}[Resistance to Reward Hacking]
\label{prop:reward_hack}
Under Pre-Squashing with scale $c$, the maximum possible advantage is $|A_i| \leq 2c$, and the escort-weighted field satisfies $|\tilde{\omega}_\alpha(A_i)| \leq \Omega_{\max}$, where $\Omega_{\max} = 2c[(1-\alpha) + \alpha e^{2c}]$. After centering, $|g_i| \leq 2\Omega_{\max}$. The BHP target satisfies:
\begin{equation}
    -1 \leq v_i^* \leq \frac{g_{\max} - g_{\min}}{\mu} + G
\end{equation}
\end{proposition}
\begin{proof}
By definition, $v_i^* = \max(-1, (g_i - \lambda^*)/\mu)$, so $v_i^* \geq -1$. For the upper bound, the bisection bracket guarantees $\lambda^* \geq g_{\min} - \mu G$ (\cref{alg:bisection}), hence $(g_i - \lambda^*)/\mu \leq (g_{\max} - g_{\min})/\mu + G$. Since $g_{\max} - g_{\min} \leq 2\Omega_{\max} \leq 4c[(1-\alpha) + \alpha e^{2c}]$, this provides a hard cap on the policy update for any single response, independent of raw reward magnitude.
\end{proof}

%====================================================================%
\section{Geometric Interpretation}
\label{sec:geometry}

\cref{fig:geometry} illustrates the geometric relationship between GOPO and W-GOPO in the Hilbert space $L^2(\pik)$.

\begin{figure}[htbp]
\centering
\tdplotsetmaincoords{65}{50}

\begin{tikzpicture}[
    tdplot_main_coords,
    scale=3.0,
    >=stealth,
    label_font/.style={font=\footnotesize},
    important_label/.style={font=\small\bfseries},
    callout/.style={thin, gray!50}
]

    \coordinate (O) at (0,0,0);

    % =================================================================================
    % 1. Background: plane H_0 (z=0) and feasible polytope K
    % =================================================================================

    \fill[gray!10, opacity=0.45] (-2.2,-2.2,0) -- (2.2,-2.2,0) -- (2.2,2.2,0) -- (-2.2,2.2,0) -- cycle;
    \draw[gray!30, thin] (-2.2,-2.2,0) -- (2.2,-2.2,0) -- (2.2,2.2,0) -- (-2.2,2.2,0) -- cycle;

    \coordinate (K1) at (0.8, -0.2, 0);
    \coordinate (K2) at (0.3, 0.8, 0);
    \coordinate (K3) at (-0.7, 0.5, 0);
    \coordinate (K4) at (-0.5, -0.6, 0);
    \coordinate (K5) at (0.5, -0.7, 0);

    \fill[green!25, opacity=0.6] (K1) -- (K2) -- (K3) -- (K4) -- (K5) -- cycle;
    \draw[green!60!black, thick] (K1) -- (K2) -- (K3) -- (K4) -- (K5) -- cycle;

    % =================================================================================
    % 2. Normal vector 1 (z-axis)
    % =================================================================================

    \draw[->, gray!55, dashed, thin] (O) -- (0,0,1.8)
        node[above, black, label_font] {$\mathbf{1}$ (Constant Function)};
    \filldraw[black] (O) circle (0.8pt);

    % =================================================================================
    % 3. GOPO pathway: u*_gopo (centered) -> proj -> outside K (conservation broken)
    % =================================================================================

    \coordinate (U_gopo) at (-0.6, 0.9, 0.6);
    \coordinate (U_gopo_plane) at (-0.6, 0.9, 0);
    \coordinate (V_gopo_bad) at (-0.6, 0.9, 0);

    \draw[->, thick, red!45, dashed] (O) -- (U_gopo);
    \draw[dashed, red!25, thin] (U_gopo) -- (U_gopo_plane);
    \filldraw[red!40] (U_gopo_plane) circle (0.5pt);

    % =================================================================================
    % 4. W-GOPO pathway: u*_wgopo (escort) -> exact lambda* -> BHP -> v* in K
    % =================================================================================

    \coordinate (U_wgopo) at (0.1, 0.8, 1.4);
    \coordinate (U_wgopo_proj) at (0.1, 0.8, 0);
    \coordinate (V_wgopo) at (0.05, 0.65, 0);

    % Vertical projection (chemical potential subtraction)
    \draw[dashed, red!50, line width=0.8pt] (U_wgopo) -- (U_wgopo_proj);

    % BHP truncation arrow
    \draw[->, dotted, orange!80!black, thick, line width=1.2pt] (U_wgopo_proj) -- (V_wgopo);
    \filldraw[red!60] (U_wgopo_proj) circle (0.6pt);

    % Main vectors
    \draw[->, very thick, red, line cap=round] (O) -- (U_wgopo);
    \draw[->, very thick, blue, line cap=round] (O) -- (V_wgopo);

    % =================================================================================
    % 5. Labels
    % =================================================================================

    % GOPO target label
    \node[label_font, red!45, align=left, anchor=south east, xshift=-2pt] at (U_gopo) {
        GOPO: $u^*_{\text{centered}}$\\[-1pt]
        {\tiny ($\lambda^*\!=\!0$ assumed)}
    };

    % GOPO landing point label (outside K, conservation broken)
    \node[label_font, red!55, anchor=north, yshift=-6pt, align=center] at (V_gopo_bad) {
        {\tiny $\sum v_i^* > 0$}\\[-1pt]
        {\tiny (conservation broken)}
    };

    % W-GOPO target label
    \node[label_font, red, align=left, anchor=south west, xshift=3pt] at (U_wgopo) {
        W-GOPO Target\\[-1pt]
        $u^*_{g} = g/\mu$
    };

    % Chemical potential label
    \node[label_font, red!65!black, anchor=west, xshift=5pt] at ($(U_wgopo)!0.55!(U_wgopo_proj)$) {
        $\lambda^*_{\text{bisect}}$ (Exact)
    };

    % BHP truncation label
    \node[label_font, orange!85!black, anchor=west, xshift=5pt, yshift=3pt] at ($(U_wgopo_proj)!0.5!(V_wgopo)$) {
        BHP Truncation
    };

    % W-GOPO policy label
    \node[important_label, blue, align=center, anchor=north east, xshift=-3pt, yshift=-22pt] at (V_wgopo) {
        W-GOPO Policy\\$v^*\!=\!P_{\mathcal{K}}(u^*_{g})$
    };

    % Reference policy label
    \node[label_font, anchor=south east, xshift=-3pt, yshift=4pt] at (O) {Reference $\pik$ ($v\!=\!0$)};

    % Feasible polytope label
    \node[label_font, green!40!black, anchor=east, align=right, xshift=-8pt] at (K5) {
        Feasible Polytope\\$\mathcal{K} = \mathcal{H}_0 \cap \mathcal{C}_+$
    };

    % H_0 label
    \node[label_font, gray!60!black, align=center,
          fill=white, fill opacity=0.85, text opacity=1,
          inner sep=3pt, rounded corners=1pt] at (1.3, 1.55, 0) {
        $\mathcal{H}_0$: Probability Conserving Subspace\\
        $\langle v, \mathbf{1} \rangle_{\pi_k} = 0$
    };

\end{tikzpicture}

\caption{\textbf{Geometric interpretation of W-GOPO vs.\ GOPO in $L^2(\pik)$.} The reference policy $\pik$ sits at the origin ($v=0$). Valid policies must reside in $\mathcal{H}_0$ (gray plane) and satisfy non-negativity, restricting them to the feasible polytope $\mathcal{K}$ (green). \emph{GOPO} (dashed red): the centered target $u^*_{\text{centered}}$ is projected with $\lambda^*=0$ assumed; when BHP truncation is active, the result may land outside $\mathcal{K}$ with $\sum v_i^* > 0$ (conservation broken). \emph{W-GOPO} (solid red/blue): the escort-weighted target $u^*_{g}$ (from centered escort field $g$) is first projected vertically by the exact chemical potential $\lambda^*_{\text{bisect}}$ (computed via \cref{alg:bisection}), then truncated onto $\mathcal{K}$ via BHP, yielding the final policy $v^*$ (blue) with strict probability conservation guaranteed.}
\label{fig:geometry}
\end{figure}

%====================================================================%
\section{Related Work}
\label{sec:related}

\paragraph{Hilbert Space Alignment.}
OPO~\cite{Wang2025OPO} first proposed lifting alignment into $L^2(\pik)$, deriving the optimal update via the Hilbert Projection Theorem and introducing the $\alpha$-Escort sampling field. GOPO~\cite{Wang2025GOPO} specialized this to finite-dimensional group subspaces, proving the vanishing chemical potential under group centering. W-GOPO is the natural synthesis: it inherits OPO's escort field and GOPO's group structure, adding exact conservation enforcement and robust reward handling.

\paragraph{KL-Based Methods.}
PPO~\cite{Schulman2017PPO} uses ratio clipping and KL penalty; DPO~\cite{Rafailov2023DPO} derives a closed-form KL-constrained solution; GRPO~\cite{Shao2024GRPO} eliminates the critic via group-relative advantages; DAPO~\cite{Yu2025DAPO} introduces asymmetric clip bounds. All inherit the exponential geometry of KL divergence, leading to gradient saturation in high-confidence regimes---the fundamental limitation addressed by the Hilbert space framework.

\paragraph{Reward Robustness.}
Reward overoptimization~\cite{Gao2023Scaling} occurs when policies exploit reward model imperfections. W-GOPO's $\tanh$ Pre-Squashing provides a defense complementary to reward model ensembles~\cite{Coste2024Reward} and conservative optimization.

\paragraph{$f$-Divergence Trust Regions.}
The $\chi^2$ divergence arising from the Hilbert space framework belongs to the $f$-divergence family~\cite{Csiszar1967,AliSilvey1966}. Unlike KL, $\chi^2$ induces constant curvature, a property that extends naturally to the W-GOPO setting.

%====================================================================%
\section{Experiments}
\label{sec:experiments}

We evaluate W-GOPO against four baselines under identical experimental conditions from the GOPO evaluation~\cite{Wang2025GOPO}, ensuring fair comparison.

\subsection{Setup}

\paragraph{Model and Data.} Qwen3-1.7B trained with the VERL framework~\cite{Sheng2024HybridFlow} on 4$\times$ RTX 4090 GPUs. Training set: $\sim$10\% of MATH Level 3 (sampled with seed 42). Validation: 100 MATH Level~4 problems---a strictly harder difficulty tier not seen during training---evaluated every 10 steps to measure out-of-distribution generalization. Shared hyperparameters: batch size 48, learning rate $2 \times 10^{-6}$, 8 epochs, $G=6$ rollouts per prompt.

\paragraph{Baselines.}
\begin{itemize}[leftmargin=2em]
    \item \textbf{OPO}~\cite{Wang2025OPO}: Full Hilbert projection with importance-sampling reweighting.
    \item \textbf{GOPO}~\cite{Wang2025GOPO}: Group Hilbert projection with group centering, $\lambda^* = 0$.
    \item \textbf{GRPO}~\cite{Shao2024GRPO}: Token-level policy gradient with ratio clipping ($\epsilon = 0.2$).
    \item \textbf{DAPO}~\cite{Yu2025DAPO}: Asymmetric clip bounds with entropy management.
\end{itemize}

\paragraph{W-GOPO Configuration.} Stiffness $\mu = 0.5$, escort exponent $\alpha = 0.5$, squash scale $c = 1.0$, bisection tolerance $\epsilon = 10^{-5}$, max iterations $T = 40$.

\subsection{Results}

\begin{table}[htbp]
\centering
\caption{Comparison on MATH benchmarks. All methods use Qwen3-1.7B trained on $\sim$10\% MATH Level~3, validated on 100 MATH Level~4 problems. Mean Reward is averaged over the entire training process. $\lambda^*$ Exact indicates whether the chemical potential is computed exactly (via bisection or global projection) or approximated ($\lambda^* = 0$ under $z$-score).}
\label{tab:results}
\begin{tabular}{lccccc}
\toprule
\textbf{Algorithm} & \textbf{Mean Reward}$\uparrow$ & \textbf{Val Acc (L4)}$\uparrow$ & \textbf{Grad Norm} & \textbf{Entropy} & \textbf{$\lambda^*$} \\
\midrule
GRPO & 0.544 & 44\% & 0.674 & 0.115 & N/A \\
DAPO & 0.548 & 44\% & 0.213 & 0.126 & N/A \\
OPO  & 0.558 & 48\% & 1.279 & 0.126 & Exact \\
GOPO & 0.555 & 47\% & 1.029 & 0.134 & Approx. \\
\textbf{W-GOPO (Ours)} & \textbf{0.561} & \textbf{49\%} & 1.152 & \textbf{0.138} & Exact \\
\bottomrule
\end{tabular}
\end{table}

\paragraph{Analysis.}
\begin{itemize}[leftmargin=2em]
    \item \textbf{Generalization.} W-GOPO achieves the highest validation accuracy (49\%) on the strictly harder MATH Level~4 problems, surpassing OPO (48\%), GOPO (47\%), and clipping-based methods (44\%). This confirms that the combination of escort weighting and exact chemical potential yields more effective learning signals than either approach alone.
    \item \textbf{Training Efficiency.} W-GOPO's 0.561 mean reward over the training process exceeds both OPO (0.558) and GOPO (0.555), indicating more efficient utilization of training data.
    \item \textbf{Gradient Dynamics.} W-GOPO maintains a healthy gradient norm (1.152) between OPO's 1.279 and GOPO's 1.029, consistent with the escort field modulating update magnitudes based on absolute batch quality. Both GRPO (0.674) and DAPO (0.213) show substantially lower gradient norms, confirming the gradient saturation endemic to clipping-based methods.
    \item \textbf{Entropy Preservation.} W-GOPO preserves the highest entropy (0.138) among all methods, indicating that the escort weighting effectively suppresses premature mode collapse by attenuating updates for degenerate batches.
\end{itemize}

%====================================================================%
\section{Discussion}
\label{sec:discussion}

\paragraph{The OPO $\to$ W-GOPO $\to$ GOPO Hierarchy.}
The three Hilbert-space algorithms form a natural theoretical hierarchy. OPO operates in the full infinite-dimensional $L^2(\pik)$ with exact global projections---theoretically complete but computationally demanding. GOPO collapses to group subspaces with maximal simplification ($\lambda^* = 0$, no escort)---practically elegant but structurally limited. W-GOPO occupies the principled middle ground: it retains GOPO's finite-dimensional group structure and critic-free design while recovering OPO's absolute quality awareness and exact conservation, adding only two scalar hyperparameters ($\alpha$, $c$) and a negligible-cost bisection subroutine.

\paragraph{When Does W-GOPO Matter Most?}
The escort weighting and exact $\lambda^*$ become critical when: (i) training data quality is heterogeneous (some prompts generate mostly poor responses); (ii) the reward model has calibration noise or adversarial outliers; (iii) BHP truncation is frequently active (small $\mu$ or large advantage magnitudes); (iv) long training runs accumulate gradual drift from degenerate batches. For short runs with high-quality data and conservative $\mu$, GOPO remains a reasonable baseline.

\paragraph{Limitations.}
W-GOPO introduces two additional hyperparameters ($\alpha$, $c$) beyond GOPO's $\mu$. While the bisection solver adds negligible wall-clock overhead, it increases implementation complexity. Our experiments focus on mathematical reasoning; validation on diverse alignment tasks (instruction following, safety, code generation) remains future work. Systematic ablation of the interaction between $\alpha$, $c$, and $\mu$ would further clarify the design space.

%====================================================================%
\section{Conclusion}
\label{sec:conclusion}

We have identified two structural pathologies of GOPO: reward amplification in degenerate batches due to group-relative centering, and probability conservation breakdown under Bounded Hilbert Projection truncation. To address these, we presented \textbf{W-GOPO}---a strict generalization of GOPO that integrates $\alpha$-Escort weighted advantages for absolute quality awareness, GPU-parallel bisection for exact chemical potential computation, and Pre-Squashing with Post-Projection for robust reward handling. The resulting algorithm preserves the Hilbert projection framework's defining properties---constant curvature, non-saturating gradients, and exact sparsity---while achieving the highest out-of-distribution generalization among all tested methods on mathematical reasoning benchmarks. W-GOPO demonstrates that the full power of the Hilbert Projection Theorem can be harnessed in practice without sacrificing the engineering simplicity that makes group-based alignment attractive.

%====================================================================%
\begin{thebibliography}{99}\setlength{\itemsep}{0pt}

\bibitem{Christiano2017RLHF}
P.~Christiano, J.~Leike, T.~B.~Brown, M.~Martic, S.~Legg, and D.~Amodei.
Deep Reinforcement Learning from Human Preferences.
\emph{NeurIPS}, 2017.

\bibitem{Ouyang2022InstructGPT}
L.~Ouyang, J.~Wu, X.~Jiang, D.~Almeida, C.~L.~Wainwright, P.~Mishkin, et al.
Training Language Models to Follow Instructions with Human Feedback.
\emph{NeurIPS}, 2022.

\bibitem{Schulman2017PPO}
J.~Schulman, F.~Wolski, P.~Dhariwal, A.~Radford, and O.~Klimov.
Proximal Policy Optimization Algorithms.
arXiv:1707.06347, 2017.

\bibitem{Rafailov2023DPO}
R.~Rafailov, A.~Sharma, E.~Mitchell, S.~Ermon, C.~D.~Manning, and C.~Finn.
Direct Preference Optimization: Your Language Model is Secretly a Reward Model.
\emph{NeurIPS}, 2023.

\bibitem{Shao2024GRPO}
Z.~Shao, P.~Wang, Q.~Zhu, R.~Xu, J.~Song, X.~Bi, H.~Zhang, M.~Zhang, Y.~K.~Li, Y.~Wu, and D.~Guo.
DeepSeekMath: Pushing the Limits of Mathematical Reasoning in Open Language Models.
arXiv:2402.03300, 2024.

\bibitem{Wang2025OPO}
Z.~Wang.
Orthogonalized Policy Optimization: Policy Optimization as Orthogonal Projection in Hilbert Space.
arXiv:2601.12415, 2026.

\bibitem{Wang2025GOPO}
Z.~Wang.
Group Orthogonalized Policy Optimization: Group Policy Optimization as Orthogonal Projection in Hilbert Space.
arXiv:2602.21269, 2026.

\bibitem{Yu2025DAPO}
Q.~Yu \emph{et al.}
DAPO: An Open-Source LLM Reinforcement Learning System.
arXiv:2503.14476, 2025.

\bibitem{Sheng2024HybridFlow}
G.~Sheng, C.~Xu, S.~Yuan, et al.
HybridFlow: A Flexible and Efficient RLHF Framework.
arXiv:2409.19256, 2024.

\bibitem{Gao2023Scaling}
L.~Gao, J.~Schulman, and J.~Hilton.
Scaling Laws for Reward Model Overoptimization.
\emph{ICML}, 2023.

\bibitem{Coste2024Reward}
T.~Coste, U.~Anwar, R.~Kirk, and D.~Krueger.
Reward Model Ensembles Help Mitigate Overoptimization.
\emph{ICLR}, 2024.

\bibitem{Csiszar1967}
I.~Csisz\'ar.
Information-type measures of difference of probability distributions and indirect observations.
\emph{Studia Sci.\ Math.\ Hungar.}, 2:299--318, 1967.

\bibitem{AliSilvey1966}
S.~M.~Ali and S.~D.~Silvey.
A General Class of Coefficients of Divergence of One Distribution from Another.
\emph{JRSS-B}, 28(1):131--142, 1966.

\end{thebibliography}

%====================================================================%
\appendix
\section{Proof of Conservation Breakdown}
\label{app:conservation}

\begin{proof}[Proof of \cref{prop:conservation_break}]
Let $S = \{i : A_i/\mu < -1\}$ with $|S| = k > 0$. Then:
\begin{align}
    \sum_{i=1}^G v_i^* &= \sum_{i \notin S} \frac{A_i}{\mu} + \sum_{i \in S} (-1) \\
    &= \frac{1}{\mu}\sum_{i=1}^G A_i - \frac{1}{\mu}\sum_{i \in S} A_i + \sum_{i \in S}(-1) \\
    &= 0 - \sum_{i \in S}\left(\frac{A_i}{\mu} + 1\right) > 0
\end{align}
since $A_i/\mu + 1 < 0$ for each $i \in S$.
\end{proof}

\section{Convergence of Bisection}
\label{app:bisection}

\begin{proposition}[Convergence Rate]
The bisection algorithm reduces the bracket width $|\lambda_{\text{high}} - \lambda_{\text{low}}|$ below $\epsilon$ in at most
\begin{equation}
    T = \left\lceil \log_2\!\left(\frac{g_{\max} - g_{\min} + \mu(G+1)}{\epsilon}\right) \right\rceil
\end{equation}
iterations, where $g_{\max} = \max_i g_i$ and $g_{\min} = \min_i g_i$ are the extrema of the centered escort field. Here $\epsilon$ is the $\lambda$-precision (bracket width); the conservation residual at the returned midpoint satisfies $|f(\lambda)| \leq (G/\mu)\cdot\epsilon/2$, since $f$ has Lipschitz constant at most $G/\mu$ (each of the $G$ unclamped summands contributes slope $1/\mu$).
\end{proposition}

\begin{proof}
The initial interval width is $(g_{\max} + \mu) - (g_{\min} - \mu G) = g_{\max} - g_{\min} + \mu(G+1)$. Bisection halves the interval at each step, yielding width $< \epsilon$ after $T$ iterations. The residual bound follows from $|f(\lambda) - f(\lambda^*)| = |f(\lambda)| \leq (G/\mu)|\lambda - \lambda^*| \leq (G/\mu) \cdot \epsilon/2$.
\end{proof}

\section{Proof of Gradient Suppression}
\label{app:suppression}

\begin{proof}[Proof of \cref{thm:suppression}]
By definition~\eqref{eq:escort}, $|\tilde{\omega}_\alpha(A_i)| = |(1-\alpha)A_i + \alpha A_i e^{A_i}| \leq (1-\alpha)|A_i| + \alpha |A_i| e^{|A_i|}$. Therefore
\begin{equation}
    \sum_{i=1}^G |\tilde{\omega}_\alpha(A_i)| \leq (1-\alpha)\sum|A_i| + \alpha \sum |A_i| e^{|A_i|}
\end{equation}
The first term equals $(1-\alpha)G\bar{A}$, scaling linearly in $\bar{A}$. For the second term, note that $h(x) = x e^x$ is convex on $[0,\infty)$, so by Jensen's inequality $\frac{1}{G}\sum h(|A_i|) \geq h(\bar{A})$, and the bound in \cref{thm:suppression} follows. The key observation is that for small $\bar{A}$, $|A_i| e^{|A_i|} \approx |A_i|(1 + |A_i|) = |A_i| + A_i^2$, so the second term behaves as $O(G\bar{A} + G\bar{A}^2)$, providing superlinear ($O(\bar{A}^2)$) suppression in the escort-active component.
\end{proof}

\end{document}
